\documentclass[11pt,a4paper]{article}
\usepackage[T1]{fontenc}
\usepackage[utf8]{inputenc}
\usepackage{lmodern}
\usepackage[margin=26mm,headheight=14pt]{geometry}
\usepackage{amsmath,amssymb,amsthm,mathtools}
\usepackage{booktabs,longtable,array}
\usepackage{microtype}
\usepackage{xcolor}
\usepackage{fancyhdr}
\usepackage{enumitem}
\usepackage[hyphens]{url}
\usepackage{xurl}
\usepackage[colorlinks=true,linkcolor=blue!45!black,citecolor=blue!45!black,urlcolor=blue!45!black]{hyperref}
\usepackage{bookmark}
\hypersetup{pdftitle={HODGE_SIGN: Sign-only purity and the Hodge-index comparison},pdfauthor={HODGE_SIGN research run, 21 November 2026}}
\setlength{\parindent}{0pt}
\setlength{\parskip}{5pt plus 1pt minus 1pt}
\setlength{\emergencystretch}{2em}
\setlist{topsep=4pt,itemsep=2pt}
\pagestyle{fancy}
\fancyhf{}
\fancyhead[L]{\small HODGE\_SIGN}
\fancyhead[R]{\small Checked mathematical report}
\fancyfoot[C]{\thepage}
\newtheorem{theorem}{Theorem}[section]
\newtheorem{lemma}[theorem]{Lemma}
\newtheorem{proposition}[theorem]{Proposition}
\newtheorem{corollary}[theorem]{Corollary}
\theoremstyle{remark}
\newtheorem{remark}[theorem]{Remark}
\DeclareMathOperator{\Tr}{Tr}
\DeclareMathOperator{\CH}{CH}
\DeclareMathOperator{\NS}{NS}
\DeclareMathOperator{\Spec}{Spec}
\DeclareMathOperator{\End}{End}
\DeclareMathOperator{\Hom}{Hom}
\title{\textbf{HODGE\_SIGN}\\[5pt]
\large Sign-only purity, evanescent cohomology,\\
and the limits of the Hodge-index comparison}
\author{Research run HODGE\_SIGN--20261121\\
\small GPT-6 Astra Pro, root and three independent task agents}
\date{21 November 2026}
\begin{document}
\maketitle
\begin{abstract}
We reconstruct the sign-only mechanism as a complete theorem: a generic
Fourier-compatible hypersurface of an affine power has a Gysin remainder
whose signed virtual class is effective and contains a tensor power of the
lowest compactly supported cohomology.  Parseval and tensor amplification
give the expected weight bound; stratification and duality give the usual
finite-field absolute values.  The sign is cohomological parity, not
positive definiteness of an intersection form.  We compute the split
quadric and CM elliptic-power cases, track every supplied SA06 C.1 row,
give a counterexample to the displayed higher-cycle height formula in
cote 44, prove corrected versions of the openness remark in cote 34, and
separate the actual theta-lift divisibility bounds from integral Lefschetz
and Wilson matrices.  The latter comparison includes the normalized
universal ring, exact small-power Smith forms, and explicit lattice
counterexamples.  No general standard conjecture or unspecified higher
arithmetic index theorem is claimed.
\end{abstract}
\tableofcontents
\newpage
\part{Checked findings and explicitly cited background}
